% =====================================================================
\begin{thebibliography}{99}
\addcontentsline{toc}{chapter}{Bibliografía}
\markboth{Bibliografía}{}
\small

\bibitem{bv}
S.~Boyd y L.~Vandenberghe.
\newblock \emph{Convex Optimization}.
\newblock Cambridge University Press, 2004.
\newblock Cap.~3, «Convex functions», y cap.~4, «Convex optimization problems». Texto del curso EE364a de Stanford.
\newblock \url{https://web.stanford.edu/~boyd/cvxbook/}

\bibitem{ee364a}
S.~Boyd.
\newblock EE364a: Convex Optimization I.
\newblock Stanford University. Transparencias de la clase «Convex functions».
\newblock \url{https://web.stanford.edu/class/ee364a/lectures/functions.pdf}

\bibitem{cs224n}
C.~Manning \emph{et al.}
\newblock CS224N: Natural Language Processing with Deep Learning.
\newblock Stanford University, invierno de 2021. Tarea~2, «word2vec».
\newblock \url{https://web.stanford.edu/class/archive/cs/cs224n/cs224n.1214/assignments/a2.pdf}

\bibitem{cs231n}
A.~Karpathy y equipo docente de CS231n.
\newblock CS231n: Deep Learning for Computer Vision. Notas del curso, «Neural Networks Part 3: Learning and Evaluation», sección «Gradient checks».
\newblock Stanford University.
\newblock \url{https://cs231n.github.io/neural-networks-3/}

\bibitem{dauphin}
Y.~N. Dauphin, R.~Pascanu, C.~Gulcehre, K.~Cho, S.~Ganguli e Y.~Bengio.
\newblock Identifying and attacking the saddle point problem in high-dimensional non-convex optimization.
\newblock En \emph{Advances in Neural Information Processing Systems 27 (NeurIPS)}, 2014.
\newblock \url{https://arxiv.org/abs/1406.2572}

\bibitem{mml}
M.~P. Deisenroth, A.~A. Faisal y C.~S. Ong.
\newblock \emph{Mathematics for Machine Learning}.
\newblock Cambridge University Press, 2020.
\newblock Cap.~5, «Vector Calculus».
\newblock \url{https://mml-book.github.io/}

\bibitem{gbc}
I.~Goodfellow, Y.~Bengio y A.~Courville.
\newblock \emph{Deep Learning}.
\newblock MIT Press, 2016.
\newblock Cap.~4, «Numerical Computation», y §6.5, «Back-Propagation and Other Differentiation Algorithms».
\newblock \url{https://www.deeplearningbook.org/}

\bibitem{mit1802}
D.~Auroux.
\newblock 18.02: Multivariable Calculus.
\newblock Massachusetts Institute of Technology, MIT OpenCourseWare, primavera de 2006.
\newblock \url{https://ocw.mit.edu/courses/18-02-multivariable-calculus-spring-2006/}

\bibitem{cookbook}
K.~B. Petersen y M.~S. Pedersen.
\newblock \emph{The Matrix Cookbook}.
\newblock Technical University of Denmark, versión del 15 de noviembre de 2012.

\bibitem{d2l}
A.~Zhang, Z.~C. Lipton, M.~Li y A.~J. Smola.
\newblock \emph{Dive into Deep Learning}, versión 1.0.3.
\newblock Cambridge University Press, 2023.
\newblock §2.4, «Calculus».
\newblock \url{https://d2l.ai/}

\end{thebibliography}
